\chapter{Analysis of the Harder--Narasimhan filtration}
\label{chap:hn}

% archon:covers MR4736527GeometricLocalSystemsOnVeryGeneralCurvesAndIsomonodromy/HNFiltration.lean

This chapter carries out the analysis of \S6 of Landesman--Litt: the numerical
constraints on the Harder--Narasimhan filtration of an isomonodromic deformation
to a general nearby curve, culminating in \cref{cor:stable-parabolic}, the
parabolic semistability statement on which the Hodge-theoretic results depend.

\section{Generic global generation}

\begin{definition}[Generic global generation]
  \label{def:generic-global-generation}
  \lean{LandesmanLitt.GenericallyGloballyGenerated}
  % SOURCE: [Landesman--Litt], \S6, Definition (generic global generation) (read from references/MR4736527-local-systems-isomonodromy.tex)
  % SOURCE QUOTE: "A vector bundle $V$ is {\em generically globally generated} if the
  % evaluation map $H^0(C, V) \otimes \mathscr O_C \to V$ does not factor through a proper
  % subbundle of $V$, i.e.~if the cokernel of this map is torsion."
  \textit{Source: Landesman--Litt, \S6.}
  A vector bundle $V$ on a smooth proper curve $C$ is \emph{generically globally
  generated} if the evaluation map
  \[
    H^0(C, V) \otimes \sO_C \longrightarrow V
  \]
  does not factor through a proper subbundle of $V$; equivalently, if the cokernel
  of this map is a torsion sheaf, i.e.\ the evaluation map is surjective at the
  generic point of $C$. A parabolic sheaf $E_\star$ is called generically globally
  generated if its underlying vector bundle $E = E_0$ is.
\end{definition}

\section{Preliminary bounds from Clifford's theorem}

\begin{theorem}[Clifford's theorem for vector bundles]
  \label{thm:clifford-vector-bundles}
  \lean{LandesmanLitt.clifford_vb}
  % SOURCE: [Brambila-Paz--Grzegorczyk--Newstead], Theorem 2.1, as cited in [Landesman--Litt], \S6 proof of lemma:cohomology-bound (read from references/MR4736527-local-systems-isomonodromy.tex)
  % SOURCE QUOTE: "Using Clifford's theorem for vector bundles
  % \cite[Theorem 2.1]{brambila-pazGN:geography-of-brill-noether-loci},
  % for $i > k$,
  % we have
  % \begin{align*}
  % \dim H^0(C, W_i)
  % \leq \frac{\deg W_i}{2} + \rk W_i.
  % \end{align*}"
  \textit{Source: Brambila-Paz--Grzegorczyk--Newstead, Theorem 2.1 (cited in Landesman--Litt, \S6).}
  Let $V$ be a semistable vector bundle on a smooth proper curve $C$ of genus $g$
  whose slope satisfies $0 \leq \mu(V) \leq 2g - 2$. Then
  \[
    \dim H^0(C, V) \leq \frac{\deg V}{2} + \rk V.
  \]
\end{theorem}

% SOURCE QUOTE PROOF: "L.~Brambila-Paz, I.~Grzegorczyk, and P.~E. Newstead.
% \newblock Geography of {B}rill-{N}oether loci for small slopes.
% \newblock {\em J. Algebraic Geom.}, 6(4):645--669, 1997."  (read from references/MR4736527-local-systems-isomonodromy-source/isomonodromy-stability-hodge-theory.bbl)
\begin{proof}
  This is the Clifford-type bound for semistable vector bundles on curves,
  Theorem~2.1 of Brambila-Paz, Grzegorczyk and Newstead, \emph{Geography of
  Brill--Noether loci for small slopes} (J.\ Algebraic Geom.\ \textbf{6} (1997),
  645--669). We take it as an external input; no Mathlib analogue is available.
\end{proof}

\begin{lemma}[Cohomology bound for slopes in $[0,2g]$]
  \label{lem:cohomology-bound}
  \lean{LandesmanLitt.cohomology_bound}
  \uses{thm:clifford-vector-bundles, def:parabolic-HN}
  % SOURCE: [Landesman--Litt], \S6, lemma:cohomology-bound (read from references/MR4736527-local-systems-isomonodromy.tex)
  % SOURCE QUOTE: "Suppose $V$ is a vector bundle on a smooth proper curve $C$ with Harder-Narasimhan filtration $0 = N^0 \subset N^1 \subset \cdots \subset N^m = V$. Suppose moreover that for each $i$, the slope of $\on{gr}^i_{N} V=N^i/N^{i-1}$
  % satisfies $$0\leq \mu(\on{gr}^i_{N} V):=\frac{\deg(\on{gr}^i_{N}V)}{\on{rk}(\on{gr}^i_{N}V)}\leq 2g.$$
  % Then $\dim H^0(C, V) \leq \frac{\deg V}{2} + \rk V$."
  \textit{Source: Landesman--Litt, \S6.}
  Let $V$ be a vector bundle on a smooth proper curve $C$ of genus $g$ with
  Harder--Narasimhan filtration $0 = N^0 \subset N^1 \subset \cdots \subset N^m =
  V$, and suppose that for each $i$ the slope of $\gr^i_N V = N^i / N^{i-1}$
  satisfies
  \[
    0 \leq \mu(\gr^i_N V) := \frac{\deg(\gr^i_N V)}{\rk(\gr^i_N V)} \leq 2g.
  \]
  Then
  \[
    \dim H^0(C, V) \leq \frac{\deg V}{2} + \rk V.
  \]
\end{lemma}

% SOURCE QUOTE PROOF: "For convenience set $W_i := \on{gr}^i_{N}V=N^i/N^{i-1}$.
% Suppose $W_1, \ldots, W_k$ have slopes $> 2g-2$ and $W_{k+1}, \ldots,
% W_m$ have slopes $\leq 2g-2$.
% Using Clifford's theorem for vector bundles
% \cite[Theorem 2.1]{brambila-pazGN:geography-of-brill-noether-loci},
% for $i > k$, we have $\dim H^0(C, W_i) \leq \frac{\deg W_i}{2} + \rk W_i.$
% Also, for $i \leq k$, since $W_i$ are semistable, there are no maps $W_i \to
% \omega_C$. Therefore, $H^1(C, W_i) = 0$ when $i \leq k$. It follows from Riemann
% Roch that $\dim H^0(C, W_i) = \deg W_i + (1-g) \rk W_i$ for $i \leq k$. Summing
% over $i$, we get $\dim H^0(C, W) \leq \frac{\deg W}{2} + \rk W+ \sum_{i=1}^k
% (\frac{\deg W_i}{2} -g \rk W_i).$ To conclude, it is enough to show $\frac{\deg
% W_i}{2} -g \rk W_i \le 0$. However, since we were assuming the slope $\mu(W_i)
% \leq 2g$, we find $\deg W_i \leq 2g \rk W_i$ and so $\frac{\deg W_i}{2} \leq g
% \rk W_i$, as desired."  (read from references/MR4736527-local-systems-isomonodromy.tex)
\begin{proof}
  \uses{thm:clifford-vector-bundles}
  Write $W_i := \gr^i_N V = N^i / N^{i-1}$, and split the indices: say
  $W_1, \ldots, W_k$ have slopes $> 2g - 2$ and $W_{k+1}, \ldots, W_m$ have slopes
  $\leq 2g - 2$. Each $W_i$ is semistable, being a Harder--Narasimhan graded piece.

  For $i > k$ the hypothesis gives $0 \leq \mu(W_i) \leq 2g - 2$, so Clifford's
  theorem for vector bundles (\cref{thm:clifford-vector-bundles}) yields
  $\dim H^0(C, W_i) \leq \tfrac{\deg W_i}{2} + \rk W_i$. For $i \leq k$ the slope
  exceeds $2g - 2$, so a semistable $W_i$ admits no nonzero map to $\omega_C$;
  hence $H^1(C, W_i) = 0$ and Riemann--Roch gives
  $\dim H^0(C, W_i) = \deg W_i + (1-g)\rk W_i$.

  Summing $\dim H^0(C, V) \leq \sum_i \dim H^0(C, W_i)$ and regrouping,
  \[
    \dim H^0(C, V) \leq \frac{\deg V}{2} + \rk V
      + \sum_{i=1}^k \Big( \frac{\deg W_i}{2} - g\,\rk W_i \Big).
  \]
  Each summand is $\leq 0$: the hypothesis $\mu(W_i) \leq 2g$ gives
  $\deg W_i \leq 2g\,\rk W_i$, hence $\tfrac{\deg W_i}{2} \leq g\,\rk W_i$. The
  stated bound follows.
\end{proof}

\section{The Harder--Narasimhan constraint}

\begin{theorem}[Parabolic Harder--Narasimhan constraints]
  \label{thm:hn-constraints-parabolic}
  \lean{LandesmanLitt.hn_constraints_parabolic}
  \uses{lem:cohomology-bound, def:generic-global-generation, def:parabolic-HN,
        prop:connection-h1, lem:parabolic-extension-obstruction,
        prop:irreducibility-splitting}
  % SOURCE: [Landesman--Litt], \S6, theorem:hn-constraints-parabolic (read from references/MR4736527-local-systems-isomonodromy.tex)
  % SOURCE QUOTE: "Let $(C,D)$ be hyperbolic of genus $g$ with $D = x_1 + \cdots + x_n$, and let
  % $({E}, \nabla)$ be a flat
  % vector bundle on $C$ with regular singularities along $D$ and irreducible
  % monodromy. Let $E_\star$ denote a parabolic structure on $E$ refined by the parabolic structure on $E$ associated to $\nabla$.
  % Suppose $(E_\star',\nabla')$ is an isomonodromic deformation (which exists by \autoref{remark:refinement})
  % of $({E_\star}, \nabla)$ to an analytically general nearby curve,
  % with Harder-Narasimhan filtration $0 = (F_\star')^0 \subset (F_\star')^1 \subset \cdots \subset
  % (F_\star')^m = E_\star'$. For $1 \leq i \leq m$, let $\mu_i$ denote the slope of
  % $\on{gr}^{i}_{HN}E_\star' := (F'_\star)^i/(F_\star')^{i-1}$.
  % Then the following two properties hold.
  % \begin{enumerate}
  %   \item If $E_\star'$ is not parabolically semistable, then for every $0 < i < m$, there
  %   exists $j < i < k$ with $$\rk \on{gr}^{j+1}_{HN}E_\star'\cdot \rk
  %   \on{gr}^k_{HN}E_\star'\geq g+1.$$
  %   \item We have $0<\mu_i-\mu_{i+1}\leq 1$ for all $i<m$.
  % \end{enumerate}"
  \textit{Source: Landesman--Litt, \S6, Theorem (parabolic Harder--Narasimhan constraints).}
  Let $(C, D)$ be hyperbolic of genus $g$ with $D = x_1 + \cdots + x_n$, and let
  $(E, \nabla)$ be a flat vector bundle on $C$ with regular singularities along
  $D$ and irreducible monodromy. Let $E_\star$ be a parabolic structure on $E$
  refined by the parabolic structure associated to $\nabla$. Suppose
  $(E_\star', \nabla')$ is an isomonodromic deformation of $(E_\star, \nabla)$ to
  an analytically general nearby curve, with Harder--Narasimhan filtration
  \[
    0 = (F_\star')^0 \subset (F_\star')^1 \subset \cdots \subset (F_\star')^m =
    E_\star',
  \]
  and for $1 \leq i \leq m$ let $\mu_i$ denote the parabolic slope of
  $\gr^i_{\HN} E_\star' := (F_\star')^i / (F_\star')^{i-1}$. Then:
  \begin{enumerate}
    \item If $E_\star'$ is not parabolically semistable, then for every
      $0 < i < m$ there exist $j < i < k$ with
      \[
        \rk \gr^{j+1}_{\HN} E_\star' \cdot \rk \gr^k_{\HN} E_\star' \geq g + 1.
      \]
    \item $0 < \mu_i - \mu_{i+1} \leq 1$ for all $i < m$.
  \end{enumerate}
\end{theorem}

% SOURCE QUOTE PROOF: "The basic idea of the proof of \autoref{theorem:hn-constraints-parabolic} will be to show that any counterexample to
% \autoref{theorem:hn-constraints-parabolic} will produce a certain semistable parabolic vector bundle of high
% slope which is not generically globally generated.
% In order to see why this failure of generic global generation
% leads to a contradiction, we will need some facts about (generic) global generation of vector bundles on curves, arising from Clifford's theorem."  (read from references/MR4736527-local-systems-isomonodromy.tex)
\begin{proof}
  \uses{lem:cohomology-bound, def:generic-global-generation,
        prop:connection-h1, lem:parabolic-extension-obstruction,
        prop:irreducibility-splitting, def:parabolic-HN}
  We sketch the argument; the full proof is assembled from the deformation-theoretic
  lemmas of the chapters on parabolic bundles and isomonodromic deformations.

  The central idea is that any counterexample produces a semistable parabolic
  vector bundle of high slope that fails to be generically globally generated,
  which contradicts the cohomology bound. In detail, after deforming to an
  analytically general nearby curve one may assume the Harder--Narasimhan
  filtration of $E_\star'$ extends to a first-order neighborhood. For part (1),
  if $E_\star'$ is not parabolically semistable then for each $0 < i < m$ the
  splitting of the Atiyah sequence determined by the connection
  (\cref{prop:connection-h1}) induces, for some $j < i$ and $k \geq i+1$, a
  nonzero map $T_{C'}(-D') \to \sheafhom(\gr^{j+1}_{\HN} E_\star',
  \gr^k_{\HN} E_\star')_\star$ whose persistence under the deformation is governed
  by the parabolic extension obstruction (\cref{lem:parabolic-extension-obstruction});
  the vanishing of the induced map on $H^1$ forces the dual twist
  $\sheafhom(\gr^{j+1}_{\HN} E_\star', \gr^k_{\HN} E_\star')^\vee \otimes
  \omega_C(D)$ to fail generic global generation
  (\cref{def:generic-global-generation}). A bundle of this slope that is not
  generically globally generated has, by the section bound built from
  \cref{lem:cohomology-bound}, a rank constraint
  $\rk \gr^{j+1}_{\HN} E_\star' \cdot \rk \gr^k_{\HN} E_\star' \geq g + 1$,
  giving (1). For part (2), the same dual twist must have slope at most
  $2g - 1 + n$; since the relevant $\sheafhom$ has negative parabolic slope by the
  definition of the Harder--Narasimhan filtration (\cref{def:parabolic-HN}), the
  parabolic slope of $\sheafhom(\gr^{j+1}_{\HN} E_\star', \gr^k_{\HN} E_\star')$
  lies in $[-1, 0)$, whence $0 < \mu_{j+1} - \mu_k \leq 1$, and the intermediate
  slopes are squeezed to give $0 < \mu_i - \mu_{i+1} \leq 1$. Irreducibility of
  the monodromy (\cref{prop:irreducibility-splitting}) is what guarantees the
  destabilizing data behaves as required. The detailed multi-lemma argument is
  deferred.
\end{proof}

\begin{theorem}[Harder--Narasimhan constraints]
  \label{thm:hn-constraints}
  \lean{LandesmanLitt.hn_constraints}
  \uses{thm:hn-constraints-parabolic}
  % SOURCE: [Landesman--Litt], \S1, theorem:hn-constraints (read from references/MR4736527-local-systems-isomonodromy.tex)
  % SOURCE QUOTE: "Let $(C,D)$ be hyperbolic of genus $g$ and let
  % $({E}, \nabla)$ be a flat
  % vector bundle on $C$ with regular singularities along $D$,
  % and irreducible
  % monodromy.
  % Suppose $(E',\nabla')$ is an isomonodromic deformation
  % of $({E}, \nabla)$ to an analytically general nearby curve,
  % with Harder-Narasimhan filtration $0 = (F')^0 \subset (F')^1 \subset \cdots \subset
  % (F')^m = E'$. For $1 \leq i \leq m$, let $\mu_i$ denote the slope of
  % $\on{gr}^{i}_{HN}E' := (F')^i/(F')^{i-1}$.
  % Then the following two properties hold.
  % \begin{enumerate}
  %   \item If $E'$ is not semistable, then for every $0 < i < m$, there
  %   exists $j < i < k$ with $$\rk \on{gr}^{j+1}_{HN}E'\cdot \rk
  %   \on{gr}^k_{HN}E'\geq g+1.$$
  %   \item We have $0<\mu_i-\mu_{i+1}\leq 1$ for all $i<m$.
  % \end{enumerate}"
  \textit{Source: Landesman--Litt, \S1, Theorem (Harder--Narasimhan constraints).}
  Let $(C, D)$ be hyperbolic of genus $g$ and let $(E, \nabla)$ be a flat vector
  bundle on $C$ with regular singularities along $D$ and irreducible monodromy.
  Suppose $(E', \nabla')$ is an isomonodromic deformation of $(E, \nabla)$ to an
  analytically general nearby curve, with Harder--Narasimhan filtration
  $0 = (F')^0 \subset (F')^1 \subset \cdots \subset (F')^m = E'$, and for
  $1 \leq i \leq m$ let $\mu_i$ denote the slope of
  $\gr^i_{\HN} E' := (F')^i / (F')^{i-1}$. Then:
  \begin{enumerate}
    \item If $E'$ is not semistable, then for every $0 < i < m$ there exist
      $j < i < k$ with
      $\rk \gr^{j+1}_{\HN} E' \cdot \rk \gr^k_{\HN} E' \geq g + 1$.
    \item $0 < \mu_i - \mu_{i+1} \leq 1$ for all $i < m$.
  \end{enumerate}
\end{theorem}

% SOURCE QUOTE PROOF: "\autoref{theorem:hn-constraints} is a special case of
% \autoref{theorem:hn-constraints-parabolic} below, where we allow certain
% parabolic structures on the vector bundle $E$."  (read from references/MR4736527-local-systems-isomonodromy.tex)
\begin{proof}
  \uses{thm:hn-constraints-parabolic}
  This is the special case of \cref{thm:hn-constraints-parabolic} in which $E$
  carries the trivial parabolic structure $E_\star = E_0$ with respect to the
  empty divisor (equivalently, unipotent monodromy at infinity, so the parabolic
  weights vanish). Under that specialization parabolic semistability is ordinary
  semistability and the parabolic slopes $\mu_i$ are ordinary slopes, so both
  conclusions follow directly.
\end{proof}

\section{Low-rank stability}

\begin{corollary}[Parabolic stability in low rank]
  \label{cor:stable-parabolic}
  \lean{LandesmanLitt.stable_parabolic}
  \uses{thm:hn-constraints-parabolic}
  % SOURCE: [Landesman--Litt], \S6, cor:stable-parabolic (read from references/MR4736527-local-systems-isomonodromy.tex)
  % SOURCE QUOTE: "Let $(C, D)$ be a hyperbolic curve of genus $g$.
  % Let $({E}, \nabla)$ be a flat
  % vector bundle on $C$ with irreducible monodromy and with regular singularities
  % along $D$, and suppose that
  % $\on{rk}(E)<2\sqrt{g+1}$.
  % Then, for $E_\star$ any parabolic structure refined by the parabolic structure on $E$ associated to $\nabla$,
  % the isomonodromic deformation of $E_\star$
  % to an analytically general nearby curve is parabolically semistable."
  \textit{Source: Landesman--Litt, \S6, Corollary (parabolic stability in low rank).}
  Let $(C, D)$ be a hyperbolic curve of genus $g$, and let $(E, \nabla)$ be a flat
  vector bundle on $C$ with irreducible monodromy and regular singularities along
  $D$, with $\rk(E) < 2\sqrt{g+1}$. Then for any parabolic structure $E_\star$
  refined by the parabolic structure associated to $\nabla$, the isomonodromic
  deformation of $E_\star$ to an analytically general nearby curve is parabolically
  semistable.
\end{corollary}

% SOURCE QUOTE PROOF: "If $(E_\star', \nabla')$ is an isomonodromic deformation of $(E_\star,\nabla)$ to an
% analytically general nearby curve which is not semistable, it follows from
% \autoref{theorem:hn-constraints-parabolic}(1) that for each $i$, there will be $j,k$ with
% $j<i<k$ so that the Harder-Narasimhan filtration $HN$ of $E_\star'$ satisfies
% $\rk \on{gr}^{j+1}_{HN}E_\star'\cdot \rk \on{gr}^k_{HN}E_\star'\geq g+1.$
% Since $\rk \on{gr}^{j+1}_{HN}E_\star' + \rk \on{gr}^k_{HN}E_\star' \leq \rk E_\star' = \rk
% E_\star$,
% it follows from the arithmetic mean-geometric mean inequality that
% $$g+1\leq \rk \on{gr}^{j+1}_{HN}E_\star'\cdot \rk \on{gr}^k_{HN}E_\star'\leq \left(\frac{\rk
% E_\star}{2} \right)^2.$$
% So $\rk E_\star \geq 2 \sqrt{g+1}$ as desired."  (read from references/MR4736527-local-systems-isomonodromy.tex)
\begin{proof}
  \uses{thm:hn-constraints-parabolic}
  Suppose the isomonodromic deformation $(E_\star', \nabla')$ were not parabolically
  semistable. Then by \cref{thm:hn-constraints-parabolic}(1) there are indices
  $j < i < k$ with
  $\rk \gr^{j+1}_{\HN} E_\star' \cdot \rk \gr^k_{\HN} E_\star' \geq g + 1$. Since
  these two graded pieces are disjoint sub quotients,
  $\rk \gr^{j+1}_{\HN} E_\star' + \rk \gr^k_{\HN} E_\star' \leq \rk E_\star' =
  \rk E_\star$, so the arithmetic mean--geometric mean inequality gives
  \[
    g + 1 \leq \rk \gr^{j+1}_{\HN} E_\star' \cdot \rk \gr^k_{\HN} E_\star'
      \leq \Big( \frac{\rk E_\star}{2} \Big)^2,
  \]
  forcing $\rk E_\star \geq 2\sqrt{g+1}$, contrary to hypothesis. Hence the
  deformation is parabolically semistable.
\end{proof}

\begin{corollary}[Stability in low rank]
  \label{cor:stable}
  \lean{LandesmanLitt.stable}
  \uses{cor:stable-parabolic}
  % SOURCE: [Landesman--Litt], \S1, cor:stable (read from references/MR4736527-local-systems-isomonodromy.tex)
  % SOURCE QUOTE: "Let $(C, D)$ be a hyperbolic curve of genus $g$.
  % Let $({E}, \nabla)$ be an irreducible flat
  % vector bundle on $C$ with regular singularities along $D$, and suppose that
  % $\on{rk}(E)<2\sqrt{g+1}$.
  % Then an isomonodromic deformation of $E$
  % to an analytically general nearby curve is  semistable."
  \textit{Source: Landesman--Litt, \S1, Corollary (stability in low rank).}
  Let $(C, D)$ be a hyperbolic curve of genus $g$, and let $(E, \nabla)$ be an
  irreducible flat vector bundle on $C$ with regular singularities along $D$ and
  $\rk(E) < 2\sqrt{g+1}$. Then an isomonodromic deformation of $E$ to an
  analytically general nearby curve is semistable.
\end{corollary}

% SOURCE QUOTE PROOF: "The following corollary can be deduced directly from
% \autoref{theorem:hn-constraints-parabolic} and AM-GM.
% It is also a special case of \autoref{cor:stable-parabolic}."  (read from references/MR4736527-local-systems-isomonodromy.tex)
\begin{proof}
  \uses{cor:stable-parabolic}
  This is the special case of \cref{cor:stable-parabolic} in which $E$ carries the
  trivial parabolic structure $E_\star = E_0$ on the empty divisor, for which
  parabolic semistability is ordinary semistability.
\end{proof}
